%% kapa: exam=1 round=29 year=2018 period=2 form=A subject=accounting subject_ko=회계학 seq=17 number=57
%% source: Q-Net 감정평가사 기출문제 (cst00309) artlSeq=5202808 file=extracted/1008/5202808/2018년 제29회 감정평가사 1차시험 기출문제/2018년 제29회 감정평가사 1차시험 기출문제/2교시 A형.pdf
%% conversion: kordoc 4.24.51 -> tools/gen1.py (auto); review=proofread-sonnet; flags=none
\begin{question}{57}
(주)감평은 20x1년 초에 부여일로부터 3년의 지속적인 용역제공을 조건으로 직원 100명에게 주식선택권을 1인당 10개씩 부여하였다. 20x1년 초 주식선택권의 단위당 공정가치는 ￦150이며, 주식선택권은 20x4년 초부터 행사할 수 있다. (주)감평의 연도별 실제 퇴직자 수 및 추가퇴직 예상자 수는 다음과 같다.
\par{\small\centering
\begin{tabular}{|c|c|c|}
\hline
 & 실제 퇴직자 수 & 추가퇴직 예상자 수 \\ \hline
20x1년 말 & 5명 & 15명 \\ \hline
20x2년 말 & 8명 & 17명 \\ \hline
\end{tabular}\par}
(주)감평은 20x1년 말에 주식선택권의 행사가격을 높이는 조건변경을 하였으며, 이러한 조건변경으로 주식선택권의 단위당 공정가치가 ￦30 감소하였다. 20x2년도 인식할 보상비용은? \par
\begin{choices}
  \item ￦16,000
  \item ￦30,000
  \item ￦40,000
  \item ￦56,000
  \item ￦70,000
\end{choices}
\end{question}
\answer{2}
